% Generated from the official 2026 MDPI ACS LaTeX template.
% Phase 24 author-review build for Mathematics (MDPI).
\documentclass[mathematics,article,submit,moreauthors]{Definitions/mdpi}


\firstpage{1}
\makeatletter
\setcounter{page}{\@firstpage}
\makeatother
\pubvolume{1}
\issuenum{1}
\articlenumber{0}
\pubyear{2026}
\copyrightyear{2026}
\datereceived{ }
\daterevised{ }
\dateaccepted{ }
\datepublished{ }

\Title{Sensitivity of Two-Agent \texorpdfstring{EF$^{21}$}{EF21} to Heterogeneous Regularity under Compression}

% Final ORCID identifiers should be inserted before submission if applicable.
\Author{Pack Kwan Low $^{1}$ and Fu-Hsing Wang $^{1,*}$}
\AuthorNames{Pack Kwan Low and Fu-Hsing Wang}

\address{$^{1}$ \quad Department of Information Management,
Chinese Culture University,
55 Hwa-Kang Rd.,
Yang-Ming-Shan,
Taipei 11114,
Taiwan}

\corres{Correspondence: wfx2@ulive.pccu.edu.tw}

\abstract{
Error Feedback 21 (EF21) provides a communication-efficient framework for distributed optimization under contractive compression, but the influence of heterogeneous local regularity remains difficult to interpret directly from existing convergence expressions. For the two-agent case with heterogeneous regularity parameters, Empirical Law 4.3 of Berg Thomsen, Taylor, and Dieuleveut characterizes the predicted optimal convergence rate of EF21 as the largest real root of a cubic polynomial. In this study, we refer to this source relation as the Heterogeneity-Aware Convergence Law (HAC Law). Before conducting the sensitivity analysis, we independently reproduced this polynomial characterization and evaluated it across 10,500 parameter configurations. The reproduction confirmed numerical consistency with the stated polynomial and with the homogeneous theoretical limit on the tested domain. We then study how communication compression interacts with heterogeneous local regularity under controlled parameterizations. An equal-smoothness baseline first varies local strong-convexity imbalance while keeping average conditioning fixed. Greater imbalance reduces the available contraction margin, and stronger compression increases the relative effect. We next allow local smoothness and strong convexity to vary independently while preserving their arithmetic means. Under this parameterization, the empirical step size and one coefficient of the cubic remain invariant, while the remaining regularity dependence is captured by a single nonnegative mismatch quantity equal to the weighted variance of local condition-shape coordinates. When smoothness and strong-convexity heterogeneity are proportionally aligned, this mismatch vanishes and the cubic reduces to the homogeneous controlled case. Misalignment makes the mismatch positive and increases the predicted contraction factor. Symbolic analysis further shows that the cubic has three distinct roots between zero and one and that its selected largest root increases with the coefficient carrying the mismatch term. These results distinguish heterogeneity magnitude from regularity mismatch and provide a structural interpretation of heterogeneous EF21 behavior under compression. All conclusions remain conditional on HAC Law and do not constitute a general convergence theorem.
}
\keyword{distributed optimization; error feedback; \texorpdfstring{EF$^{21}$}{EF21}; communication compression; regularity heterogeneity; performance estimation; sensitivity analysis} 

% Submission-draft display guard: the MDPI class synthesizes a placeholder DOI
% from volume/issue/article-number defaults even in submit mode. Suppress that
% footer in the author-facing draft; the Editorial Office assigns the real DOI.
\makeatletter
\AtBeginDocument{%
  \ifthenelse{\equal{\@status}{submit}}{%
    \rfoot{}%
    \fancypagestyle{plain}{\fancyfoot[R]{}}%
  }{}%
}
\makeatother

\setlength{\headheight}{19pt}

\begin{document}

% Section files are split for repository maintainability.
\section{Introduction}

The rapid growth of distributed and federated learning has significantly increased the importance of communication-efficient optimization algorithms. In modern large-scale learning systems, a large number of agents cooperatively train a global model while communicating with a central server. Since the transmitted gradient vectors often contain millions or even billions of parameters, communication remains one of the primary bottlenecks in large-scale distributed and federated optimization, motivating the development of gradient quantization, sparsification, and other compressed communication techniques. While these methods substantially reduce communication overhead, they inevitably introduce compression errors that may deteriorate convergence performance. To mitigate this issue, error-feedback mechanisms accumulate and re-inject compression residuals into subsequent iterations, thereby compensating for information loss and recovering favorable convergence properties even under biased or lossy compression schemes \citep{seide2014onebit,stich2018sparsified,karimireddy2019error}. Such mechanisms have become a central component of modern communication-efficient optimization algorithms.

Among existing error-compensation approaches, \emph{Error Feedback 21} (EF$^{21}$), proposed by Richt\'arik et al., has emerged as a particularly influential framework \citep{richtarik2021ef21}. By maintaining local gradient estimators and transmitting compressed gradient differences rather than compressed gradients directly, EF$^{21}$ provides a simpler and more stable formulation of error feedback in distributed environments. Recent theoretical developments have established tight convergence guarantees, optimal step-size choices, and Lyapunov-based analyses for EF$^{21}$ under contractive compression and strongly convex objectives, demonstrating its effectiveness in communication-constrained distributed learning \citep{thomsen2026tight}.

Despite these advances, the role of agent heterogeneity remains insufficiently understood. In practical federated learning systems, local objectives often exhibit substantial differences in data distribution, smoothness, and conditioning. Although recent work has derived empirical convergence laws for heterogeneous EF$^{21}$ systems \citep{thomsen2026tight}, a systematic understanding of how heterogeneity interacts with compression error to influence convergence speed and stability is still lacking. In particular, it remains unclear whether increasing heterogeneity merely slows convergence or fundamentally changes the sensitivity of EF$^{21}$ to compression-induced perturbations.

Recent evidence suggests that the interaction between communication compression
and agent heterogeneity plays a fundamental role in determining the performance
of error-feedback algorithms. In particular, Berg Thomsen et al.
\citep{thomsen2026tight} reported empirical laws of EF and
EF$^{21}$ under heterogeneous regularity parameters. The empirical laws were not obtained through a direct theoretical derivation. Instead, they emerged from an iterative discovery process combining \emph{Performance Estimation Problem (PEP)} analyses \citep{drori2014performance,taylor2017exact}, semidefinite-programming-based Lyapunov searches, symbolic algebra exploration, and extensive numerical experimentation. We also independently reconstructed the cubic polynomial specified in the empirical law and evaluated it across the complete 10,500-configuration reproduction grid, whose parameter construction is detailed in the Methods section. The audit checks the polynomial residual in every selected root, compares the implementation with a source helper pinned to a specific commit, tests the homogeneous limit against the proven EF$^{21}$ rate, and applies mutated-polynomial negative controls. High-precision SDPA calculations provide an additional solver-level reference in selected cases. The maximum polynomial residual on the reproduction grid is $1.742\times10^{-15}$. These checks support numerical consistency with the stated polynomial on the tested domain. The resulting empirical characterization was subsequently examined through large-scale parameter sweeps covering diverse regularity configurations and compression regimes. These results provide computational evidence supporting the proposed relation on the tested heterogeneous settings. In this paper, we refer to the two-agent EF$^{21}$ characterization in Empirical Law 4.3 as the \emph{Heterogeneity-Aware Convergence Law}, hereafter abbreviated as \emph{HAC Law}, because it captures how compression effects and disparities in local objective regularity jointly shape the predicted convergence behavior across agents. More specifically, for the heterogeneous two-agent EF$^{21}$ setting, HAC Law characterizes the predicted optimal convergence rate as the largest real root of an agent-dependent cubic polynomial. This characterization shows that a global condition number alone does not fully describe the predicted contraction behavior when local smoothness and strong-convexity parameters differ. 

 

The source relation is supported by numerical Performance Estimation Problem
experiments and Lyapunov feasibility analysis. For HAC Law, the
candidate rate $\rho^\star$ is taken as the largest real root of the cubic
polynomial and the step size is set to the empirical optimum $\eta^\star$.
Feasibility of $(\rho^\star,\eta^\star)$ is checked for EF$^{21}$ and EF
using the corresponding Lyapunov linear matrix inequalities. If the EF$^{21}$
feasibility check fails, the source protocol compares $\rho^\star$ with a
simplified Lyapunov rate computed by PEPit \citep{goujaud2024pepit} and accepts the fallback check when
the difference is within $10^{-4}$. Tightness is then challenged with the
improved target $0.99\rho^\star$. The source study reports no successful
strict improvement in the verification experiments \citep{thomsen2026tight}.

Accordingly, this study quantifies the sensitivity of EF$^{21}$ to heterogeneous local regularity and compression using HAC Law as a numerical and analytical reference.

The literature uses the term ``heterogeneity'' for several distinct phenomena.
\emph{Statistical heterogeneity} describes differences in local data
distributions across workers. These distributional differences can lead to
differences in the corresponding stochastic gradients.
\emph{Regularity heterogeneity} describes differences in worker-specific
smoothness constants and strong-convexity constants. Such differences also
induce variation in local conditioning.

The distinction is important because methods designed to accommodate
statistical heterogeneity do not necessarily model a separate
strong-convexity constant for each worker. For example, later EF$^{21}$ work
improves dependence on heterogeneous smoothness constants and introduces
smoothness-aware weighting \citep{richtarik2024reloaded}. EControl also permits
worker-dependent smoothness constants, while strong convexity is imposed on
the averaged objective through a single global parameter
\citep{gao2024econtrol}.

This study focuses on worker-specific regularity heterogeneity in the paired
local parameters $(L_i,\mu_i)$. Statistical heterogeneity lies outside the
scope of the analysis. We further examine whether the smoothness and
strong-convexity ratios are proportionally aligned across the two workers.
This alignment determines the regularity-mismatch term in HAC Law,
so the effect cannot be characterized by heterogeneity magnitude alone.

Berg Thomsen, Taylor, and Dieuleveut developed tight analyses of classic error feedback and EF$^{21}$ and explicitly distinguished proved results from empirical laws in more challenging heterogeneous regimes \citep{thomsen2026tight}. The present analysis uses HAC Law for $n=2$, which allows the local smoothness and strong-convexity parameters to differ across agents. For a specified compression level and pair of local regularity parameters, the law defines an empirical optimal step size and a cubic polynomial. The largest real root of that cubic gives the predicted optimal contraction factor.

Ref.~\cite{thomsen2026tight} labels this relation an \emph{empirical law} and does not state it as a general theorem for the full heterogeneous setting. This distinction is central to the analysis. The present study therefore separates validation from structural interpretation. The source PEP and Lyapunov evidence is reviewed first, followed by an independent numerical reproduction. The sensitivity analysis is then restricted to the cubic polynomial specified in HAC Law. Every statement below about the predicted contraction factor remains conditional on this relation.

The source analysis is closely connected to the \emph{Performance Estimation Problem (PEP)} framework, in which worst-case analysis of first-order methods can be formulated as semidefinite optimization \citep{drori2014performance,taylor2017exact}. Related PEP work exploiting agent symmetries further illustrates why two-agent reductions can be structurally informative in distributed performance analysis, although under different algorithms and assumptions \citep{colla2026symmetries}.

The present study asks two nested questions:

\begin{quote}
\textbf{How do local regularity heterogeneity and communication compression jointly
affect the contraction prediction of HAC Law when average conditioning is
held fixed?}
\end{quote}

and, after both smoothness and strong convexity are allowed to vary,

\begin{quote}
\textbf{Can the apparent two-dimensional heterogeneity dependence be reduced to a
smaller structural coordinate that explains when heterogeneity is harmful and
when it is invisible to HAC Law?}

\medskip
\noindent
Under fixed arithmetic means, the apparent two-dimensional heterogeneity
dependence collapses onto a single nonnegative mismatch coordinate, as
visualized in Figure~\ref{fig:rate-collapse}.
\end{quote}

The main contributions are summarized below.

\begin{enumerate}
\item \textbf{Independent reproduction and numerical audit of HAC Law.} The PEP and Lyapunov verification protocol reported by Berg Thomsen et al.\ is examined together with an independent replay of all 10,500 cubic configurations. The replay reaches a maximum polynomial residual of $1.742\times10^{-15}$, agrees with the homogeneous theoretical limit within $8.389\times10^{-13}$, rejects all mutated-polynomial controls, and finds no one-percent improvement in the audited high-precision reference checks.
\item \textbf{Controlled parameterization and baseline landscape.} A fixed-average equal-smoothness baseline is constructed for the cubic polynomial in HAC Law. The joint compression and heterogeneity landscape is evaluated over 216,027 configurations, and the resulting guidelines for contraction margin retention are checked with off-grid samples and boundary refinement.
\item \textbf{Mismatch factorization and alignment invariance.} The fixed-average design is extended to allow smoothness and strong-convexity heterogeneity to vary simultaneously. The gap between the source coefficients of HAC Law is shown to be exactly the weighted variance of local condition-shape coordinates. Under this parameterization, the gap factors through $(\tau_L-\tau_\mu)^2$, so proportional alignment leaves the cubic unchanged.
\item \textbf{Generic root sensitivity.} The root structure of the source cubic is examined on the stated admissible domain. The cubic has three distinct roots in $(0,1)$, its selected largest root exceeds $\sqrt{\epsilon}$, and $\mathrm{d}\rho^\star/\mathrm{d}K_1>0$. Together with the fixed-average factorization, this implies that every admissible nonzero mismatch strictly worsens the predicted contraction factor relative to the aligned path.
\end{enumerate}

\section{Background and HAC Law}

This section states HAC Law for the two-agent case used as the source relation for the present study and separates the quantities reproduced from Ref.~\cite{thomsen2026tight} from the algebraic reinterpretations developed here. HAC Law is not proved in the present work. All subsequent statements about its predicted contraction factor are conditional on the reproduced form of this source relation.

\subsection{Distributed objective and communication compression}

Consider the finite-sum distributed objective

\[
f(x)=\frac{1}{n}\sum_{i=1}^{n} f_i(x),
\]

with a central server coordinating local workers. Worker $i$ is associated with a smoothness constant $L_i$ and, in the heterogeneous strongly convex setting considered by HAC Law, a local strong-convexity constant $\mu_i$. Communication is compressed by a contractive operator with error level $\epsilon\in(0,1)$, consistent with the EF$^{21}$ setting and the reproduced source analysis \citep{richtarik2021ef21,thomsen2026tight}.

\subsection{Source relation from HAC Law}

The empirical relation reported by Berg Thomsen, Taylor, and Dieuleveut and referred to here as HAC Law addresses the heterogeneous EF$^{21}$ setting with $n=2$ \citep{thomsen2026tight}. The two workers may have different local smoothness constants $L_i$ and local strong-convexity constants $\mu_i$, while communication is governed by a contractive compression level $\epsilon\in(0,1)$.

For $n=2$, define

\[
\Sigma_i=L_i+\mu_i,\qquad
\Delta_i=L_i-\mu_i.
\]

The two regularity-dependent coefficients appearing in the source relation are

\[
K_1=
\frac{\Delta_2^2\Sigma_1+\Delta_1^2\Sigma_2}
{\Sigma_1\Sigma_2(\Sigma_1+\Sigma_2)},
\qquad
K_2=
\left(
\frac{\Delta_1+\Delta_2}
{\Sigma_1+\Sigma_2}
\right)^2.
\]

Let

\[
s=\sqrt{\epsilon},
\qquad
r(s)=\frac{(1-s)^2}{1+s}.
\]

HAC Law expresses the predicted optimal contraction factor $\rho^\star$ as the largest real root of the cubic polynomial \citep{thomsen2026tight}.

\[
Q(\rho)=\rho^3-A\rho^2+B\rho-s^4,
\]

where

\[
A=s(2+s)+r(s)(sK_1+K_2),
\]

and

\[
B=s^2\left[1+2s+r(s)(K_1+sK_2)\right].
\]

The same source relation associates this contraction prediction with the empirical optimal step size.

\[
\eta^\star=
\frac{4}{(L_1+\mu_1)+(L_2+\mu_2)}
\frac{1-s}{1+s}.
\]

Thus, for a specified pair of local regularity parameters and compression level, HAC Law maps the heterogeneous two-agent configuration to an empirical step size and a cubic whose largest real root gives the predicted optimal contraction factor. A smaller selected root corresponds to faster predicted linear contraction. A root closer to one corresponds to slower predicted contraction.

The label \emph{empirical law} follows Ref.~\cite{thomsen2026tight}. This manuscript uses the relation as the analytical starting point for the subsequent sensitivity analysis. The structural, numerical, and symbolic results in the following sections are conditional on its stated $n=2$ form.

\subsection{Evidence supporting the source relation}

The source paper supports the empirical relation referred to here as HAC Law through numerical Performance
Estimation Problem experiments and Lyapunov feasibility analysis. The candidate
rate $\rho^\star$ is taken as the largest real root of the cubic polynomial,
and the step size is set to the empirical optimum $\eta^\star$. The pair
$(\rho^\star,\eta^\star)$ is checked for feasibility in the EF$^{21}$ and
EF Lyapunov linear matrix inequalities. If the EF$^{21}$ feasibility check
fails, the source protocol compares $\rho^\star$ with a simplified Lyapunov
rate computed by PEPit and accepts the fallback check when the difference is
within $10^{-4}$ \citep{thomsen2026tight,goujaud2024pepit}. The source
protocol also tests

\[
\rho_{\mathrm{imp}}=0.99\rho^\star
\]

as a local tightness challenge. The source study reports no successful strict
improvement in the verification experiments \citep{thomsen2026tight}.

The $10^{-4}$ threshold is a fallback numerical agreement criterion, not a
global error bound for HAC Law. Likewise, the factor $0.99$ tests
whether a rate that is one percent smaller can be certified on the audited
instances; it does not establish that HAC Law is universally
accurate to one percent.

\subsection{Independent reproduction of the cubic polynomial in HAC Law}

A separate reproduction workflow was completed before the present sensitivity analysis. The workflow independently recomputed the cubic coefficients and selected roots over the complete 10,500-configuration grid used in the independent $n=2$ reproduction. For every configuration, it recorded the residual $|Q(\rho^\star)|$ and compared the independent calculation with the hash-pinned source helper. The homogeneous limit was also checked against the proved EF21 rate. A negative-control suite modified the polynomial and required the modified relation to fail the same checks. High-precision SDPA calculations were used as additional solver-level references on selected cases.

The independent workflow also retains an explicit solver-scope limitation. The proprietary MOSEK workload was not re-executed in full. Completed hash-pinned source transcripts were audited where those solver results were required, while open-solver checks were executed independently. The purpose of this validation layer is therefore to test numerical reproducibility and consistency with available PEP and Lyapunov evidence. It does not promote HAC Law to a theorem.

\subsection{Weighted-moment reinterpretation of the source coefficients}\label{sec:weighted-moment}

To expose the regularity structure of the source coefficients, define the \emph{local condition-shape coordinate}

\[
q_i=
\frac{\Delta_i}{\Sigma_i}
=
\frac{L_i-\mu_i}{L_i+\mu_i},
\]

and the corresponding normalized weights

\[
w_i=
\frac{\Sigma_i}{\Sigma_1+\Sigma_2}.
\]

With these definitions, the source coefficients can be written exactly as

\[
K_1=\sum_{i=1}^{2} w_iq_i^2,
\qquad
K_2=\left(\sum_{i=1}^{2} w_iq_i\right)^2.
\]

Here, $K_1$ is the weighted second moment of the local condition-shape coordinates. The coefficient $K_2$ is the square of their weighted mean. Consequently,

\[
K_1-K_2
=
\operatorname{Var}_w(q_i)
:=
\sum_{i=1}^{2}
w_i
\left(
q_i-\sum_{j=1}^{2}w_jq_j
\right)^2.
\]

This identity is an exact algebraic rewriting of the coefficients appearing in HAC Law; it is not a new convergence law. Its role in the present study is interpretive: it shows that the coefficient gap measures the weighted dispersion of the local condition-shape coordinates. The fixed-average parameterization introduced later uses this representation to determine when regularity heterogeneity activates a mismatch penalty.

\subsection{Nearby heterogeneity analyses}

Prior EF$^{21}$ work already provides meaningful analyses of heterogeneity along different axes. Error Feedback Reloaded improves dependence on heterogeneous smoothness constants from a quadratic mean to an arithmetic mean and introduces smoothness-aware weighting \citep{richtarik2024reloaded}. Its regularity description centers on smoothness heterogeneity and does not model the paired local $(L_i,\mu_i)$ condition-shape mismatch studied here.

EControl is another important neighboring analysis. It allows worker-dependent smoothness constants and imposes strong convexity on the averaged objective through a single global parameter; individual local objectives need not be strongly convex \citep{gao2024econtrol}. The present study addresses a complementary regime by examining how paired local $(L_i,\mu_i)$ structure interacts with compression within HAC Law. Smoothness-only and single-$\mu$ models do not represent this paired structure.


\section{Methods}

\subsection{Independent validation metrics}

The independent validation reports four complementary quantities. First, the cubic residual $|Q(\rho^\star)|$ checks whether the selected numerical root satisfies the stated polynomial. Second, the implementation difference compares the independently computed result with the pinned source helper. Third, the homogeneous-limit difference compares the heterogeneous cubic at the homogeneous parameter limit with the proved EF21 contraction rate from Theorem 3.1 of Ref.~\cite{thomsen2026tight}. Fourth, solver-level tightness checks test whether the improved target $0.99\rho^\star$ can be certified on the audited high-precision reference cells.

Negative controls are included to distinguish reproducibility from a permissive numerical test. Each of the 10,500 grid cells is paired with a mutated-polynomial control. A valid validation procedure should accept the source polynomial and reject the modified relation. The validation outcomes are reported before the sensitivity results so that the numerical status of the source relation is explicit before it is used as an analytical reference.

\subsection{Equal-smoothness baseline}

The baseline controlled path sets

\[
n=2,\qquad L_1=L_2=L=1,
\]

and fixes the average strong convexity

\[
\bar\mu=\frac{\mu_1+\mu_2}{2}.
\]

Strong-convexity heterogeneity is parameterized by

\[
\tau=\frac{\mu_2}{\mu_1},\qquad 0<\tau\le 1.
\]

Holding $\bar{\mu}$ fixed gives

\[
\mu_1=\frac{2\bar\mu}{1+\tau},
\qquad
\mu_2=\frac{2\bar\mu\tau}{1+\tau}.
\]

Thus, varying $\tau$ changes the imbalance in local strong convexity while preserving its arithmetic mean. The homogeneous case corresponds to $\tau=1$.

\subsection{Conditioning strata and baseline outcomes}

With $L=1$, define

\[
\bar\kappa=\frac{L}{\bar\mu}.
\]

The baseline studies $\bar{\kappa}\in \{2,10,100\}$. For each $(\tau,\epsilon,\bar{\kappa})$ cell, the analysis records the empirical step size, the selected largest real root $\rho^\star$, the homogeneous baseline $\rho_{\mathrm{hom}}$ from Theorem 3.1 of Ref.~\cite{thomsen2026tight}, the absolute penalty

\[
H_{\rm abs}=\rho^\star-\rho_{\mathrm{hom}},
\]

the normalized penalty

\[
H_{\rm norm}=
\frac{\rho^\star-\rho_{\mathrm{hom}}}
{1-\rho_{\mathrm{hom}}},
\]

and contraction-margin retention

\[
R=
\frac{1-\rho^\star}
{1-\rho_{\mathrm{hom}}}
=1-H_{\rm norm}.
\]

This normalization prevents a small absolute difference near $\rho=1$ from being interpreted as negligible when the remaining contraction margin is itself already small.

\subsection{Baseline dense study and retention boundaries}

The primary baseline grid uses

\begin{itemize}
\item $\tau\in [0.05,1]$, 381 evenly spaced points;
\item $\epsilon\in [0.01,0.95]$, 189 evenly spaced points;
\item $\bar{\kappa}\in \{2,10,100\}$.
\end{itemize}

This yields

\[
381\times 189\times 3=216{,}027
\]

controlled configurations.

For a target retention $R_0$, the one-sided operating boundary is the smallest studied $\tau$ satisfying

\[
R(\tau,\epsilon,\bar\kappa)\ge R_0.
\]

The main guideline uses $R_0=0.99$; finite-grid estimates are compared with bisection references at $\epsilon=0.95$.

\subsection{Baseline off-grid robustness}

The deterministic off-grid stress test evaluates 20,000 paired off-grid samples per conditioning stratum. Each comparison tests whether the numerical ordering with respect to $\tau$ or $\epsilon$ is consistent with the dense-grid pattern, while also recording the residual of the cubic at each selected root. These checks assess numerical robustness; they do not constitute proofs of global monotonicity.

\subsection{Full-regularity fixed-average parameterization}

The full extension releases the equal-smoothness constraint and fixes both arithmetic means:

\[
\bar L=\frac{L_1+L_2}{2},
\qquad
\bar\mu=\frac{\mu_1+\mu_2}{2},
\qquad
\bar\kappa=\frac{\bar L}{\bar\mu}.
\]

Smoothness and strong-convexity heterogeneity are varied independently through

\[
\tau_L=\frac{L_2}{L_1},
\qquad
\tau_\mu=\frac{\mu_2}{\mu_1},
\qquad
0<\tau_L,\tau_\mu\le 1.
\]

The controlled reconstruction is

\[
L_1=\frac{2\bar L}{1+\tau_L},
\qquad
L_2=\tau_L L_1,
\]

\[
\mu_1=\frac{2\bar\mu}{1+\tau_\mu},
\qquad
\mu_2=\tau_\mu\mu_1.
\]

A cell is admissible only when $0<\mu_i\le L_i$ for both workers. Combinations that violate these local regularity conditions are excluded from the audit; no projection or coercion is applied.

The default full-regularity audit uses 31 values for each heterogeneity ratio, 95 compression values, and three conditioning strata:

\[
31\times31\times95\times3=273{,}885
\]

requested cells, of which 258,400 satisfy the local regularity constraints.


\subsection{Fixed-average mismatch factorization}

Using the weighted-moment representation established in Section~\ref{sec:weighted-moment}, we now examine the source coefficients under the fixed-average parameterization. Because both arithmetic means are fixed, the empirical step size from HAC Law is constant across $(\tau_L,\tau_\mu)$ for fixed $(\bar L,\bar{\mu},\epsilon)$. The same controlled parameterization also yields the exact invariant

\[
K_2=\left(\frac{\bar\kappa-1}{\bar\kappa+1}\right)^2.
\]

The remaining regularity dependence is carried by $K_1-K_2$. Algebraic simplification yields

\[
K_1-K_2=
\frac{4\bar\kappa^2(\tau_L-\tau_\mu)^2}
{(\bar\kappa+1)^2
(\tau_L+\bar\kappa\tau_\mu+\bar\kappa+1)
(\bar\kappa\tau_L\tau_\mu+\tau_L\tau_\mu+\bar\kappa\tau_L+\tau_\mu)}.
\]

Every denominator factor is positive on the admissible positive domain. Hence

\[
K_1\ge K_2,
\]

with equality exactly when

\[
\tau_L=\tau_\mu.
\]

This aligned path represents proportional local regularity scaling: the workers may remain heterogeneous in their raw $L_i$ and $\mu_i$ values, while the smoothness and strong-convexity ratios vary in lockstep.

\subsection{Generic symbolic root and sensitivity audit}

The full-regularity structure motivates a cubic-coordinate analysis independent of any particular $(\tau_L,\tau_\mu)$ grid. For the studied $\bar{\kappa}>1$ setting, positivity of the local regularity parameters and the weighted-moment representation imply

\[
0<K_2\le K_1<1,
\qquad
0<s<1.
\]

A separate Wolfram Language verification script analyzes the source cubic on this generic domain. The symbolic certificate establishes a strictly positive discriminant, excludes roots at or below zero and at or above one, and evaluates

\[
Q(s)=K_2(s-1)^3s^2<0.
\]

Thus all three roots are distinct and lie in $(0,1)$, while the selected largest root satisfies

\[
\rho^\star>s=\sqrt\epsilon.
\]

Implicit differentiation with respect to $K_1$ gives

\[
\frac{d\rho^\star}{dK_1}
=
\frac{(1-s)^2}{1+s}
\frac{s\rho^\star(\rho^\star-s)}{Q'(\rho^\star)}.
\]

At the largest simple real root, $Q'(\rho^\star)>0$. Therefore

\[
\frac{d\rho^\star}{dK_1}>0.
\]

\subsection{AI-Assisted Research Workflow and Verification}
AI-assisted tools were used in study design, code development, computational planning, and manuscript preparation; the scope of this assistance is disclosed in the Acknowledgments. Symbolic claims reported in this manuscript were independently checked with Wolfram Language 15.0.1.


\section{Results}

\subsection{Independent validation outcome}

The independent replay covers all 10,500 cubic configurations. The maximum polynomial residual is $1.742\times10^{-15}$, and the maximum difference from the pinned source helper is zero at the stored numerical precision. At the homogeneous limit, the largest-root prediction agrees with the proved EF21 rate to within $8.389\times10^{-13}$. All 10,500 mutated-polynomial controls are rejected. The high-precision SDPA audit contains two reference cases and 72 associated improved-rate cells; none certifies the target $0.99\rho^\star$.

These checks address two separate questions. The residual, helper comparison, homogeneous limit, and negative controls establish reproducibility of the cubic calculation. The PEP and Lyapunov checks address whether the candidate rate is consistent with the available worst-case numerical certificates. Taken together, the source evidence and the independent reproduction provide a numerical basis for using the cubic in the subsequent sensitivity study. They do not provide a universal approximation-error bound and do not replace a proof of HAC Law.

\subsection{Equal-smoothness contraction landscape}

Figure 1 maps $\rho^\star(\tau,\epsilon)$ for $\bar{\kappa}=10$. The predicted contraction becomes slower as compression error increases, and the heterogeneous surface separates progressively from the homogeneous edge as $\tau$ decreases.

\begin{figure}[H]
\centering
\includegraphics[width=0.96\textwidth]{figures/figure1_contraction_landscape.pdf}
\caption{Optimal contraction factor across heterogeneity and compression at fixed average conditioning. The surface is evaluated from the cubic polynomial in HAC Law at $\bar{\kappa}=10$, with $\tau=\mu_2/\mu_1$ controlling worker heterogeneity while $(\mu_1+\mu_2)/2$ is held fixed. Smaller $\tau$ corresponds to stronger heterogeneity and larger $\epsilon$ to heavier compression error. The figure visualizes the joint variation of the largest admissible real root $\rho^\star$; values closer to one indicate slower linear contraction. The plot summarizes the cubic over the audited domain.\label{fig:contraction-landscape}}
\end{figure}
\unskip


\subsection{Relative heterogeneity penalty}

Figure 2 maps $H_{\mathrm{norm}}$ for the same conditioning stratum. Across the audited baseline grid, stronger strong-convexity heterogeneity increases the normalized penalty, while heavier compression amplifies the relative burden within the stated numerical tolerance.

Across the dense grid, the maximum normalized penalties are approximately:

\begin{itemize}
\item 5.07\% for $\bar{\kappa}=2$;
\item 1.70\% for $\bar{\kappa}=10$;
\item 0.20\% for $\bar{\kappa}=100$.
\end{itemize}

\begin{figure}[H]
\centering
\includegraphics[width=0.96\textwidth]{figures/figure2_normalized_penalty.pdf}
\caption{Relative loss of contraction margin caused by heterogeneity at $\bar{\kappa}=10$. The normalized heterogeneity penalty is $H_{\mathrm{norm}}=(\rho^\star-\rho_{\mathrm{hom}})/(1-\rho_{\mathrm{hom}})$, where $\rho_{\mathrm{hom}}$ is the homogeneous theoretical baseline evaluated at the same average conditioning and compression level. This normalization expresses the heterogeneous slowdown as a fraction of the contraction margin remaining under the homogeneous baseline.\label{fig:normalized-penalty}}
\end{figure}
\unskip


\subsection{Interaction between compression and conditioning}

At the strongest audited baseline heterogeneity, $\tau=0.05$, increasing $\epsilon$ from $0.01$ to $0.95$ multiplies the normalized penalty by approximately:

\begin{itemize}
\item $3.84\times$ for $\bar{\kappa}=2$;
\item $3.12\times$ for $\bar{\kappa}=10$;
\item $3.03\times$ for $\bar{\kappa}=100$.
\end{itemize}

The smaller relative penalty under poor conditioning should not be interpreted as faster convergence. In the $\bar{\kappa}=100$ dense study, every cell satisfies $\rho^\star\ge0.95$, and approximately 63.5\% satisfy $\rho^\star\ge0.99$. This stratum is already close to the non-contractive limit, leaving comparatively little contraction margin for heterogeneity to consume.

\begin{figure}[H]
\centering
\includegraphics[width=0.96\textwidth]{figures/figure3_conditioning_interaction.pdf}
\caption{Compression amplifies the relative heterogeneity burden, with the magnitude depending on baseline conditioning. Curves show the normalized heterogeneity penalty versus compression error at the strongest audited heterogeneity ($\tau=0.05$) for $\bar{\kappa}\in\{2,10,100\}$. From $\epsilon=0.01$ to $0.95$, the normalized penalty increases by approximately $3.84\times$, $3.12\times$, and $3.03\times$ for $\bar{\kappa}=2,10,100$, respectively. The smaller relative penalty at $\bar{\kappa}=100$ should not be interpreted as faster convergence because that stratum is already close to the non-contractive limit over much of the audited domain.\label{fig:conditioning-interaction}}
\end{figure}
\unskip


\subsection{Contraction-margin operating boundaries}

Figure 4 converts the normalized penalty into a 99\% retention boundary. At $\epsilon=0.95$, bisection gives

\begin{itemize}
\item $\tau^\star=0.4076217486520454$ for $\bar{\kappa}=2$;
\item $\tau^\star=0.17985615951449502$ for $\bar{\kappa}=10$;
\item full-domain satisfaction for the audited range $\tau\ge 0.05$ when $\bar{\kappa}=100$.
\end{itemize}

These thresholds provide operating summaries of the cubic within the audited domain.

\begin{figure}[H]
\centering
\includegraphics[width=0.96\textwidth]{figures/figure4_retention_boundary.pdf}
\caption{Minimum heterogeneity ratio required to retain 99\% of the homogeneous contraction margin. For each compression level and conditioning stratum, the boundary reports the smallest audited $\tau$ satisfying $R=(1-\rho^\star)/(1-\rho_{\mathrm{hom}})\ge 0.99$. At $\epsilon=0.95$, continuous bisection references are $\tau^\star=0.4076217487$ for $\bar{\kappa}=2$ and $\tau^\star=0.1798561595$ for $\bar{\kappa}=10$; the full audited range $\tau\ge0.05$ satisfies the criterion for $\bar{\kappa}=100$. These are operating boundaries for the cubic within the audited domain.\label{fig:retention-boundary}}
\end{figure}
\unskip

\begin{table}[H]
\caption{Summary of the controlled equal-smoothness EF$^{21}$ sensitivity analysis by average-conditioning stratum. The table reports the maximum normalized and absolute heterogeneity penalties, the observed range of $\rho^\star$, the 99\% retention boundary at $\epsilon=0.95$, and whether the full audited range satisfies the 99\% criterion.\label{tab:key-results}}
\centering
\begin{tabular}{cccccc}
\toprule
$\bar{\kappa}$ & Max. $H_{\mathrm{norm}}$ & Max. $H_{\mathrm{abs}}$ & $\rho^\star_{\min}$ & $\rho^\star_{\max}$ & $\tau^\star_{99\%}$\\
\midrule
2   & 0.050736 & 0.015650 & 0.250000 & 0.983908 & 0.407622\\
10  & 0.016990 & 0.002103 & 0.728505 & 0.995427 & 0.179856\\
100 & 0.002005 & 0.000030 & 0.967907 & 0.999493 & $\le 0.05$\\
\bottomrule
\end{tabular}
\end{table}


\subsection{Boundary and off-grid robustness}

Figure 5 compares finite-grid boundary estimates with bisection references as the number of $\tau$ grid points increases. The principal boundary values remain stable under refinement to 1521 grid points. The deterministic off-grid audit likewise identifies no substantive ordering violation under the project's $1\times10^{-9}$ numerical interpretation tolerance, while the maximum cubic-root residual remains approximately $2\times10^{-15}$.

\begin{figure}[H]
\centering
\includegraphics[width=0.96\textwidth]{figures/figure5_boundary_convergence.pdf}
\caption{Stability of the 99\% retention boundary under increasing $\tau$-grid resolution. Finite-grid estimates at $\epsilon=0.95$ are compared with bisection references for the three conditioning strata. The grid-derived thresholds converge toward the continuous references as resolution increases, supporting the interpretation that the operating boundaries are not artifacts of a coarse parameter grid. The bisection procedure still evaluates the cubic numerically and therefore does not constitute an analytic convergence proof.\label{fig:boundary-convergence}}
\end{figure}
\unskip


\subsection{Fixed-stratum symbolic audit}

For each $\bar{\kappa}\in \{2,10,100\}$, the discriminant can also be factored separately into a positive rational prefactor, $(1-s)^6s^4$, and a bivariate polynomial in $(s,\tau)$. All 63 coefficients of this polynomial are strictly positive. The cubic therefore has three distinct real roots throughout the corresponding open equal-smoothness domains.

This calculation is restricted to the three selected conditioning strata. It is kept as an independent check; the generic cubic-coordinate certificate given below covers a broader domain.



\subsection{Full-regularity heterogeneity reduces to a mismatch coordinate}

Allowing both $L_i$ and $\mu_i$ to vary does not change the empirical step size or $K_2$ as long as their arithmetic means are fixed. The remaining two-dimensional regularity dependence is therefore carried by the nonnegative gap $K_1-K_2$.

The weighted-variance identity gives this gap a direct meaning. It measures dispersion in the local condition-shape coordinates $q_i$ and does not separately measure variation in $L_i$ or $\mu_i$. The dense audit contains 273,885 requested configurations, of which 258,400 satisfy $\mu_i\le L_i$. On these admissible cells, the moment identity, the invariant value of $K_2$, and the fixed-average factorization agree to machine precision.

\begin{figure}[H]
\centering
\includegraphics[width=0.96\textwidth]{figures/figure6_mismatch_geometry.pdf}
\caption{Full-regularity mismatch geometry at fixed average conditioning. The heatmap shows the exact source coefficient gap $K_1-K_2$ over the two-dimensional ratio plane $(\tau_L,\tau_\mu)$ for $\bar{\kappa}=10$. The dashed diagonal $\tau_L=\tau_\mu$ is the proportional-heterogeneity path and forms an exact zero-mismatch valley. Cells violating the local regularity condition $\mu_i\le L_i$ are masked and receive no artificial values. The figure visualizes the controlled fixed-average factorization of the HAC Law coefficients.\label{fig:mismatch-geometry}}
\end{figure}
\unskip


\subsection{Proportional heterogeneity leaves the cubic polynomial unchanged}

The numerator of the fixed-average factorization contains $(\tau_L-\tau_\mu)^2$, and no other numerator factor vanishes on the positive admissible domain. Hence the zero-mismatch path is exactly $\tau_L=\tau_\mu$.

The two workers need not be identical on this path. Their smoothness and strong-convexity constants may differ in magnitude, but the two regularity pairs scale proportionally. In that case, $K_1=K_2$, the empirical step size remains fixed, and the cubic coefficients are the same as in the homogeneous controlled case with the same $\bar{\kappa}$ and $\epsilon$.

The distinction is therefore between \textbf{heterogeneity magnitude} and \textbf{regularity mismatch}. Raw heterogeneity alone does not create a penalty in the source empirical relation. The penalty appears when the smoothness and strong-convexity ratios no longer align.

\begin{figure}[H]
\centering
\includegraphics[width=0.96\textwidth]{figures/figure7_full_regularity_penalty.pdf}
\caption{Full-regularity normalized contraction penalty under heavy compression. The heatmap shows the normalized penalty $(\rho^\star-\rho_{\mathrm{hom}})/(1-\rho_{\mathrm{hom}})$ at $\bar{\kappa}=10$ and $\epsilon=0.95$ across $(\tau_L,\tau_\mu)$. The aligned diagonal $\tau_L=\tau_\mu$ forms a zero-penalty valley because the cubic reduces exactly to the homogeneous controlled case on that path. Moving away from the diagonal activates the nonnegative mismatch coordinate and worsens the reproduced largest-root contraction prediction. Invalid local regularity cells remain masked.\label{fig:full-penalty}}
\end{figure}
\unskip


\subsection{Regularity mismatch strictly worsens the largest-root prediction}

\begin{Proposition}[Conditional mismatch principle]
Under the reproduced cubic polynomial specified in HAC Law for $n=2$, fix $\bar L$, $\bar\mu$, and $\epsilon$, and parameterize local regularity by $(\tau_L,\tau_\mu)$ on the admissible domain. Then the empirical step size and $K_2$ are invariant, while the cubic depends on $(\tau_L,\tau_\mu)$ only through the nonnegative gap $K_1-K_2$. This gap vanishes if and only if $\tau_L=\tau_\mu$. Moreover, the selected largest root satisfies $\mathrm{d}\rho^\star/\mathrm{d}K_1>0$ on $0<s<1$ and $0<K_2\le K_1<1$; hence every admissible nonzero mismatch strictly increases the predicted contraction factor relative to the aligned path.
\end{Proposition}

On the cubic-coordinate domain $0<s<1$, $0<K_2\le K_1<1$, the symbolic certificate places all three roots of the source cubic in $(0,1)$ and shows that they are distinct. The largest root lies above $s$ and is simple, so implicit differentiation gives

\[
\frac{d\rho^\star}{dK_1}>0.
\]

For fixed average conditioning and compression, $K_2$ is constant and $K_1=K_2+(K_1-K_2)$. The mismatch factorization therefore gives the ordering

\[
\rho^\star(\tau_L,\tau_\mu)
\ge
\rho^\star(\tau,\tau),
\]

with strict inequality for every admissible cell satisfying $\tau_L\ne \tau_\mu$.

Conditional on HAC Law, the smallest largest-root prediction on this controlled surface is attained on the aligned path.

\begin{figure}[H]
\centering
\includegraphics[width=0.96\textwidth]{figures/figure8_rate_collapse.pdf}
\caption{Collapse of two-dimensional regularity heterogeneity onto the single mismatch coordinate. For $\epsilon=0.95$ and $\bar{\kappa}\in\{2,10,100\}$, every admissible $(\tau_L,\tau_\mu)$ pair is plotted as normalized contraction penalty versus $K_1-K_2$. Within each conditioning stratum, the two-dimensional ratio pairs collapse onto a one-dimensional curve because $K_2$ and the reproduced empirical step size are fixed, with the cubic varying only through $K_1=K_2+(K_1-K_2)$. The monotone direction is consistent with the analytic sensitivity result $\mathrm{d}\rho^\star/\mathrm{d}K_1>0$.\label{fig:rate-collapse}}
\end{figure}
\unskip



\section{Discussion}

\subsection{From heterogeneity magnitude to heterogeneity mismatch}

The equal-smoothness experiment provides an initial picture in which larger local curvature imbalance is associated with a worse contraction prediction, and stronger compression increases the relative burden. The full-regularity analysis refines this interpretation. Once smoothness and strong convexity vary independently, the cubic depends on both the extent of worker differences and the alignment of the two regularity ratios.

This distinction is visible on the path $\tau_L=\tau_\mu$. The workers may still differ in their absolute regularity scales, yet they have the same local condition-shape coordinate. In that case, the gap $K_1-K_2$ is zero and the cubic gives the homogeneous controlled prediction. Away from the diagonal, the local condition-shape coordinates separate and the coefficient gap becomes positive.

\subsection{What comes from the source and what is derived here}

Two steps in the analysis should be kept separate. The identity

\[
K_1-K_2=\operatorname{Var}_w(q_i)
\]

follows directly from rewriting the coefficients in HAC Law. It identifies $K_1$ as a weighted second moment and $K_2$ as the square of the corresponding weighted mean. The factorization in $(\bar{\kappa},\tau_L,\tau_\mu)$ is obtained only after imposing the fixed-average parameterization used in this study; this is the step that makes the alignment condition explicit.

The same division applies to the root results. The cubic itself is taken from HAC Law in Ref.~\cite{thomsen2026tight}. Its numerical use in this paper is supported by the source PEP and Lyapunov verification and by the independent reproduction reported above. The discriminant calculation, root localization, and implicit sensitivity analysis are carried out here to describe the behavior of that source relation. The validation supports this use on the audited domain without establishing the empirical relation as a general theorem.

\subsection{Algebraic generalizability beyond two agents}\label{sec:generalizability}

The weighted-moment identity is algebraically valid beyond the two-agent case. For any finite set of coordinates $q_i$ and normalized nonnegative weights $w_i$ satisfying $\sum_i w_i=1$,

\[
\sum_i w_iq_i^2-
\left(\sum_i w_iq_i\right)^2
=
\operatorname{Var}_w(q_i).
\]

A coefficient pair with the same moment structure therefore admits the same variance interpretation for $n>2$. The present results do not extend the full HAC Law analysis to arbitrary $n$, however. The identification with the source coefficient structure, the fixed-average factorization, and the cubic root-sensitivity result are all tied to the $n=2$ relation in HAC Law studied here. A separate convergence or performance-estimation analysis would be needed for more agents.

\subsection{Relation to prior \texorpdfstring{EF$^{21}$}{EF21} heterogeneity analyses}

Heterogeneous smoothness has already been studied in the EF$^{21}$ family and is known to affect complexity bounds and weighting choices \citep{richtarik2024reloaded}. The mismatch coordinate considered here addresses a narrower question: the joint dependence on local $(L_i,\mu_i)$ pairs in HAC Law for $n=2$ under a fixed-average design.

EControl considers a different regularity model \citep{gao2024econtrol}. Local $L_i$ values may vary, while strong convexity is imposed on the averaged objective through a single parameter. It therefore does not contain the same paired local regularity structure. In the related-work chain examined for this study, no theorem-level counterpart of the fixed-average factorization, alignment invariance, or cubic root-sensitivity result was identified.

\subsection{Practical interpretation}

For the equal-smoothness baseline, the retention ratio expresses how much of the homogeneous contraction margin remains after heterogeneity is introduced. The full-regularity analysis adds another distinction. Proportional scaling of local smoothness and strong convexity produces no regularity-mismatch penalty in the cubic. A difference between the two ratios makes the mismatch coordinate positive and worsens the selected root.

This result provides a rate interpretation. It does not constitute a design rule for a complete distributed learning system because communication cost, stochastic effects, and several implementation details are absent from the cubic model. Its main value is the separation of two forms of heterogeneity that a conventional parameter sweep would otherwise combine.


\section{Limitations}

\begin{enumerate}
\item HAC Law is used as a source relation from Ref.~\cite{thomsen2026tight}. The present study does not establish that law in full generality.
\item The validation establishes numerical reproducibility and consistency with the audited PEP and Lyapunov evidence. The $10^{-4}$ fallback tolerance and the one-percent tightness challenge are properties of the source verification protocol, not universal accuracy bounds for HAC Law.
\item The independent reproduction did not re-execute the complete proprietary MOSEK workload. Hash-pinned completed transcripts were audited where needed, and open-solver checks were executed independently.
\item The analysis depends on the form of HAC Law in the cited version of Ref.~\cite{thomsen2026tight}. A future revision of that expression would require the fixed-average factorization and sensitivity analysis to be checked again.
\item The full-regularity cubic analysis, its fixed-average factorization, and the root-sensitivity result are limited to two agents ($n=2$). Only the weighted-moment variance identity is algebraically more general, as discussed in Section~\ref{sec:generalizability}.
\item Numerical monotonicity in the equal-smoothness study is reported only for $\tau\in [0.05,1]$, $\epsilon\in [0.01,0.95]$, and the selected conditioning strata. The off-grid calculations provide robustness evidence but no proof of global monotonicity.
\item For the factorized discriminant calculation, the conditioning values are restricted to $\bar{\kappa}\in \{2,10,100\}$. The broader root statement comes from a separate certificate on the stated $(s,K_1,K_2)$ domain.
\item Cells that violate $\mu_i\le L_i$ are excluded from the full-regularity grid. No conclusion is extended to those invalid local regularity combinations.
\item The contraction-margin boundary summarizes rate retention; it is not an optimization of communication cost.
\item The related-work audit covers the named EF literature chain. We make no claim of universal novelty beyond this scope.
\item Large stochastic learning experiments are outside this study. The cubic characterization should not be used to infer their behavior.
\end{enumerate}


\section{Conclusion}

In this paper, the sensitivity of two-agent EF$^{21}$ to heterogeneous local regularity is studied using HAC Law reported in Ref.~\cite{thomsen2026tight}. Before the sensitivity analysis, the evidential basis of the relation is examined using the source PEP and Lyapunov verification protocol and an independent reproduction of the complete 10,500-configuration cubic grid. The replay reaches a maximum polynomial residual of $1.742\times10^{-15}$ and recovers the homogeneous theoretical limit within $8.389\times10^{-13}$. These checks support numerical use of the cubic on the audited domain while preserving its empirical status.

The equal-smoothness study shows that strong-convexity heterogeneity reduces the available contraction margin and that compression increases its relative effect. The full-regularity study gives a different view of the same problem. With the arithmetic means of smoothness and strong convexity fixed, the empirical step size and $K_2$ remain unchanged, while the remaining coefficient gap is the weighted dispersion of the local condition-shape coordinates. This gap contains the squared mismatch $(\tau_L-\tau_\mu)^2$.

When the two heterogeneity ratios are aligned, the workers may still have different local regularity scales, but the cubic is unchanged from the homogeneous controlled case. A nonzero mismatch makes the gap positive and increases the predicted contraction factor. The generic root analysis also shows that the selected largest root increases strictly with $K_1$ at fixed $K_2$. The results therefore distinguish the amount of local heterogeneity from the mismatch between the two regularity ratios.

Further work may examine whether a similar mismatch coordinate appears for $n>2$, either through performance-estimation methods or a separate convergence analysis. Stochastic large-scale experiments and the possible use of the mismatch quantity in adaptive weighting or compression-aware design also remain open directions.

\authorcontributions{Author contributions will be finalized using the CRediT taxonomy before submission. All authors have read and agreed to the published version of the manuscript.}

\funding{Funding information will be confirmed before submission.}

\institutionalreview{Not applicable. This study did not involve human participants or animals.}

\informedconsent{Not applicable.}

\dataavailability{
The analysis code, numerical summaries, symbolic-verification materials,
and reproducibility scripts supporting this study are available in the
public GitHub repository at
\href{https://github.com/Jaycee871/Sensitivity-of-TwoAgent-EF21-to-Heterogeneous-Regularity-under-Compression}
{Sensitivity of Two-Agent EF21 to Heterogeneous Regularity under Compression}.
}

\acknowledgments{During the preparation of this study and manuscript, the authors used ChatGPT (OpenAI; GPT-5.6 Sol) to assist with research-workflow design, code generation and review, computational experiment planning, interpretation checks, and manuscript drafting and editing. Wolfram Language 15.0.1 was used as an independent symbolic-computation environment for algebraic simplification and discriminant analysis. All AI-assisted outputs, code, numerical and symbolic results, citations, and manuscript text were reviewed and verified by the human authors, who take full responsibility for the content of the publication.}

\conflictsofinterest{Fu-Hsing Wang serves as a Guest Editor of the Special Issue ``Artificial Intelligence and Algorithms'' in \emph{Mathematics}. Independent editorial handling by an Editorial Board member without a conflict of interest is requested. Any additional conflicts of interest will be confirmed and declared before submission.}

\abbreviations{Abbreviations}{The following abbreviations are used in this manuscript:\\
\noindent\begin{tabular}{@{}ll}
EF & Error Feedback\\
EF21 & Error Feedback 21\\
PEP & Performance Estimation Problem\\
\end{tabular}}

\reftitle{References}
\begin{thebibliography}{99}

\bibitem{seide2014onebit}
Seide, F.; Fu, H.; Droppo, J.; Li, G.; Yu, D.
1-bit stochastic gradient descent and its application to data-parallel distributed training of speech DNNs.
In \emph{Proceedings of Interspeech 2014};
Singapore, 14--18 September 2014;
pp. 1058--1062.
\href{https://doi.org/10.21437/Interspeech.2014-274}
{doi:10.21437/Interspeech.2014-274}.

\bibitem{stich2018sparsified}
Stich, S.U.; Cordonnier, J.-B.; Jaggi, M.
Sparsified SGD with memory.
In \emph{Advances in Neural Information Processing Systems};
2018; Volume 31, pp. 4452--4463.

\bibitem{karimireddy2019error}
Karimireddy, S.P.; Rebjock, Q.; Stich, S.U.; Jaggi, M.
Error feedback fixes SignSGD and other gradient compression schemes.
In \emph{Proceedings of the 36th International Conference on Machine Learning};
PMLR, 2019; Volume 97, pp. 3252--3261.

\bibitem{richtarik2021ef21}
Richt\'arik, P.; Sokolov, I.; Fatkhullin, I.
EF21: A new, simpler, theoretically better, and practically faster error feedback.
In \emph{Advances in Neural Information Processing Systems};
2021; Volume 34, pp. 4384--4396.

\bibitem{thomsen2026tight}
Berg Thomsen, D.; Taylor, A.B.; Dieuleveut, A.
A Tight Theory of Error Feedback Algorithms in Distributed Optimization.
\emph{arXiv} \textbf{2026}, arXiv:2605.31594.
\href{https://doi.org/10.48550/arXiv.2605.31594}
{\nolinkurl{doi:10.48550/arXiv.2605.31594}}.

\bibitem{drori2014performance}
Drori, Y.; Teboulle, M.
Performance of first-order methods for smooth convex minimization: A novel approach.
\emph{Math. Program.} \textbf{2014}, \emph{145}, 451--482.
\href{https://doi.org/10.1007/s10107-013-0653-0}
{doi:10.1007/s10107-013-0653-0}.

\bibitem{taylor2017exact}
Taylor, A.B.; Hendrickx, J.M.; Glineur, F.
Exact worst-case performance of first-order methods for composite convex optimization.
\emph{SIAM J. Optim.} \textbf{2017}, \emph{27}, 1283--1313.
\href{https://doi.org/10.1137/16M108104X}
{doi:10.1137/16M108104X}.

\bibitem{goujaud2024pepit}
Goujaud, B.; Moucer, C.; Glineur, F.; Hendrickx, J.M.; Taylor, A.B.; Dieuleveut, A.
PEPit: Computer-assisted worst-case analyses of first-order optimization methods in Python.
\emph{Math. Program. Comput.} \textbf{2024}, \emph{16}, 337--367.
\href{https://doi.org/10.1007/s12532-024-00259-7}
{doi:10.1007/s12532-024-00259-7}.

\bibitem{richtarik2024reloaded}
Richt\'arik, P.; Gasanov, E.; Burlachenko, K.
Error Feedback Reloaded: From Quadratic to Arithmetic Mean of Smoothness Constants.
In \emph{Proceedings of the 12th International Conference on Learning Representations (ICLR 2024)};
2024.

\bibitem{gao2024econtrol}
Gao, Y.; Islamov, R.; Stich, S.U.
EControl: Fast Distributed Optimization with Compression and Error Control.
In \emph{Proceedings of the 12th International Conference on Learning Representations (ICLR 2024)};
2024.

\bibitem{colla2026symmetries}
Colla, S.; Hendrickx, J.M.
Exploiting Agent Symmetries for Performance Analysis of Distributed Optimization Methods.
\emph{Open J. Math. Optim.} \textbf{2026}, \emph{7}, Article 1, 1--42.
\href{https://doi.org/10.5802/ojmo.49}
{doi:10.5802/ojmo.49}.

\end{thebibliography}
\end{document}

